% HRG-C3 (#13172): pelvis frame re-expression and documented HJC relocation.
% Standalone section; the HRG reference (HRG-00, #13184) can % HRG-C3 (#13172): pelvis frame re-expression and documented HJC relocation.
% Standalone section; the HRG reference (HRG-00, #13184) can % HRG-C3 (#13172): pelvis frame re-expression and documented HJC relocation.
% Standalone section; the HRG reference (HRG-00, #13184) can % HRG-C3 (#13172): pelvis frame re-expression and documented HJC relocation.
% Standalone section; the HRG reference (HRG-00, #13184) can \input{sections/pelvis_frame}.
\section{Pelvis Frame Re-Expression and Hip Centre Relocation}
\label{sec:hrg-pelvis-frame}

\paragraph{Status.} Software correctness only; scientific qualification stays in the
design-manual pathway. No refit was run: the joint trajectory is the fitted bundle's,
so the numbers below describe the model change, not a better match. The fit's inverse
dynamics are replayed state by state (inverse kinematics plus \texttt{mj\_inverse}),
not an independent open-loop simulation.

\subsection{Problem}
\texttt{apply\_hip\_calibration} moves only the hip joints. On the best driver fit
(\texttt{full-body-v1}, 44 coordinates, AN-1 upper body) the pelvis centre of mass
(COM) and L5/S1 stayed at their design offsets from the pelvis origin, which the
builder places at the hip-joint-centre (HJC) midpoint, while the midpoint itself moved
down 78.7~mm: COM 163.7~mm and L5/S1 217.3~mm above the midpoint against 84.5~mm and
138.2~mm by design. The hip lever arms were therefore wrong.

\subsection{Operation (a): Pure Re-Expression}
Let $F$ be the pelvis body frame of the specification and $F'=FT$ the new frame, with
$T=(R,t)\in SE(3)$ its pose in $F$. The MuJoCo body frame $J$ (the frame the six root
coordinates act on) is not moved, so coordinates keep their meaning. Every quantity
stored in $F$ is re-expressed with $T^{-1}$:
\begin{align}
  p' &= R^\top (p - t), &
  c' &= R^\top (c - t),\quad I'_{c} = R^\top I_{c} R, \\
  P' &= T^{-1} P \text{ (child joint frames, named frames, closure placements)}, &
  C' &= T^{-1} C \text{ (pelvis to $J$, \texttt{child\_to\_follower})},
\end{align}
with $c$ the solid COM and $I_c$ its inertia tensor about $c$. Visual ellipsoids are
axis aligned, so a rotation is refused when the body carries one (translation only).
Markers attach to named frames and are untouched.

\paragraph{Invariance sketch.} A body's world pose is $X_J\,C^{-1}$-composed with its
content: a point $p$ in $F$ is at $X_J C^{-1} p$ in the world ($C$ the pelvis-to-$J$
transform). After the change it is $X_J C'^{-1} p' = X_J (T^{-1}C)^{-1} T^{-1}p
= X_J C^{-1} T T^{-1} p = X_J C^{-1}p$. A child joint frame obeys the same identity, so
child poses are unchanged for every coordinate vector. Mass and the world inertia
$R_w I_c R_w^\top$ are unchanged because $R_w$ absorbs $R$ and $I'_c=R^\top I_c R$.
All terms of the equations of motion are functions of these world quantities and the
unchanged coordinates, so $\mathbf q$, $\dot{\mathbf q}$, $\ddot{\mathbf q}$ and the
generalised forces (\texttt{qfrc\_inverse}) are identical. Composition holds,
$\mathrm{reexpress}(T_1)\circ\mathrm{reexpress}(T_2)=\mathrm{reexpress}(T_1T_2)$, and
$T^{-1}$ undoes $T$ to $10^{-12}$.

\paragraph{Test.} \texttt{tests/unit/motion\_matching/test\_pelvis\_frame.py}: world
landmarks (frame and contact sites), per-body COM and mass, world inertia tensor,
whole-body COM and \texttt{qfrc\_inverse}, relative to the largest entry, $\le10^{-9}$,
for a translation and a $12^\circ,-25^\circ,140^\circ$ rotation plus translation, on
\texttt{full\_body\_spec\_anthro\_driver} and \texttt{full\_body\_spec\_v1}. A mutation
test (tensor left unrotated) fails the same check, so the check is not vacuous. On the
best-fit specification (private, by path): relative differences $0$ for the translation
and $\le 10^{-15}$ (COM $2.9\times10^{-16}$, world inertia $1.0\times10^{-15}$,
\texttt{qfrc\_inverse} $8.8\times10^{-16}$, sites $4.4\times10^{-16}$~m) for the
rotation, with $\max|\mathtt{qfrc}|=615$~N\,m at a random state.

\subsection{Operation (b): Documented Physical Relocation}
Inputs: HJCs $c_R,c_L$ in the \texttt{Hip} (waist marker) frame from HRG-C2
(\texttt{hip\_centre\_map}); status must be \texttt{identified} (fails closed on
\texttt{prior\_dominated} or \texttt{degenerate} unless accepted explicitly).
With $m=\tfrac12(c_R+c_L)$:
\begin{enumerate}
  \item The pelvis frame origin moves to $m$ by operation (a), axes unchanged.
  \item Hip joint translations are set to $c_R,c_L$ ($R$ blocks, hence hip axes and
  zero-twist calibration, unchanged).
  \item The pelvis COM, the L5/S1 pivots (\texttt{TorsoInput} and the spine universal
  joint) and the pelvis visual hint keep their \emph{design offsets from $m$}
  (read from the input; $d_{COM}=(0,0,0.0845)$, $d_{L5S1}=(0,0,0.1382)$~m here).
  \item The bodies hung on L5/S1 keep their fitted pose at zero twist and spine
  angles: the joint pivot moves by $\delta=m$, the follower offset takes the same shift
  ($C_u'=F_u^{-1}\,\mathrm{Trans}(\delta)\,J_u$, likewise for the spine joint). The
  twist and spine rotation centres therefore change, which is the intended effect on
  the L5/S1 torques; trunk, arms, head, markers and total mass do not move.
\end{enumerate}
The output is a NEW document (\texttt{spec\_revision} block: issue, parent canonical
hash, centres, offsets, report) with its own \texttt{spec\_sha256}; the writer refuses
to overwrite a spec or receipt. The parent and all historical receipts are untouched.

\subsection{Independent Pelvic-Plane Reference}
Aligning the pelvis $y$ axis to the hip line makes the two HJC heights equal by
construction, so equal heights then say nothing. The reference here is the waist
marker plane (the \texttt{Hip} frame $z$ axis), and the hip line is a measured quantity
with an interval (\texttt{pelvic\_plane\_report}). On the best capture the MAP height
difference is $-20.2$ mm [$-60.2$, 20.4] and the roll $-6.2^\circ$ [$-18.2$, 6.0]: the
data do not exclude a level hip line, so the frame is \emph{not} rotated to the hip
line. Caveat: the same four markers define the plane and feed the estimate, so common-mode
marker error (13.3 mm waist residual) is shared; there is no ASIS reference in this
capture.

\subsection{Results on the Best Driver Fit}
HJC displacements (original pelvis frame): right $(0.2,\,5.6,\,11.2)$ mm, left
$(2.9,\,-25.4,\,70.6)$ mm; inter-HJC 220.3 to 186.9 mm. L5/S1 to HJC midpoint 217.3 to
138.2 mm (inside 0.13--0.15 m; this equals the design offset by construction, and the
compiled-model geometry confirms it independently). Pelvis COM to midpoint 163.7 to
84.5 mm.

\begin{center}
\begin{tabular}{lcc}
\hline
Address lever arms (horizontal, m, $x$ forward, $y$ left) & Before & After \\
\hline
Pelvis COM from right HJC & $(-0.068,-0.106)$ & $(-0.037,-0.095)$ \\
Pelvis COM from left HJC & $(-0.104,\ 0.109)$ & $(-0.041,\ 0.091)$ \\
Above-hip COM from right HJC & $(-0.232,-0.104)$ & $(-0.226,-0.098)$ \\
Above-hip COM from left HJC & $(-0.269,\ 0.111)$ & $(-0.231,\ 0.088)$ \\
Above-hip COM height over right / left HJC & 0.297 / 0.328 & 0.279 / 0.258 \\
\hline
\end{tabular}
\end{center}

\paragraph{Inverse dynamics.} 91 frames (every 20th of the 1815-frame bundle, 1 ms
sampling, acceleration by central difference of the bundle velocity; impact near frame
1323). The hip reaction wrench is the force and moment about the HJC midpoint that the
legs must apply to the pelvis, from \texttt{qfrc\_inverse} through the root Jacobian,
with the legs given negligible mass (so no leg motion or ground contact is needed). The
L5/S1 torques depend only on the bodies above.
\begin{center}
\begin{tabular}{lcc}
\hline
Peak over the slice & Before & After \\
\hline
Hip reaction force (N) & 1433.5 & 1396.3 \\
Hip reaction moment about HJC midpoint (N\,m) & 316.2 & 414.2 \\
$|\tau|$ \texttt{SpineInputX} (N\,m) & 177.1 & 178.8 \\
$|\tau|$ \texttt{SpineInputY} (N\,m) & 256.2 & 355.2 \\
$|\tau|$ \texttt{TorsoInput} (N\,m) & 222.4 & 244.7 \\
\hline
\end{tabular}
\end{center}
The changes follow from the geometry: the above-hip COM, the pelvis COM and the L5/S1
pivots changed distance to the pivots while the upper-body motion did not. They do not
mean the new torques are more accurate. Stills:
\texttt{HRG\_2026-10-17/HRG-C3/before\_after\_hip\_reaction\_l5s1.png}.

\subsection{Limitations and Failed Attempts}
\begin{itemize}
  \item The left HJC rises 70.6 mm relative to the pelvis markers, so at the fitted
  joint angles the left foot is lifted by the same amount; a contact-aware whole-body
  inverse dynamics on this trajectory would mostly measure that, so none is reported.
  The refit is a later wave.
  \item A first test compared upper-body poses at random spine angles and failed:
  the bodies keep their pose only at zero twist and spine angles, because the rotation
  centre moved. The test now zeroes those coordinates and says so.
  \item The native \texttt{full\_body\_spec\_v1} has spine pivots not colocated with
  the twist pivot, so (b) refuses it (\texttt{ValueError}); (a) works on it.
  \item The root coordinates stay at the body frame $J$; moving them to $m$ needs a state
  map (Euler recovery, velocity, acceleration) and belongs to \texttt{full-body-v2}.
  \item Only a pure translation is applied in (b) and rotations are refused with visual
  ellipsoids; pelvis width is HRG-C4.
\end{itemize}

\subsection{Reproduction}
\begin{verbatim}
# tests (red before the module existed: ImportError at collection)
PYTHONPATH=.:src python3 -m pytest -o addopts="" -q -n 4 \
    tests/unit/motion_matching/test_pelvis_frame.py
# best-fit revision and comparison (private inputs by path; outputs outside the repository)
R=$HOME/fleet24/brick-sim-s32/ct/s32-fw-drv/runs
PYTHONPATH=.:src python3 -m scripts.research.hrg_c3_pelvis_revision \
    --spec $R/driver/full_body_spec_hipcal_scaled.json --receipt $R/driver/receipt.json \
    --capture data/C3D_TA_Driver.c3d --bundle $R/bundles/driver.npz \
    --out-dir $HOME/fleet24/hrg-c3 --stride 20
\end{verbatim}
Revision \texttt{spec\_sha256} 0d08a23a\ldots{} (canonical 314a7f1d\ldots{}) from parent
\texttt{d1829abf\ldots{}}.
.
\section{Pelvis Frame Re-Expression and Hip Centre Relocation}
\label{sec:hrg-pelvis-frame}

\paragraph{Status.} Software correctness only; scientific qualification stays in the
design-manual pathway. No refit was run: the joint trajectory is the fitted bundle's,
so the numbers below describe the model change, not a better match. The fit's inverse
dynamics are replayed state by state (inverse kinematics plus \texttt{mj\_inverse}),
not an independent open-loop simulation.

\subsection{Problem}
\texttt{apply\_hip\_calibration} moves only the hip joints. On the best driver fit
(\texttt{full-body-v1}, 44 coordinates, AN-1 upper body) the pelvis centre of mass
(COM) and L5/S1 stayed at their design offsets from the pelvis origin, which the
builder places at the hip-joint-centre (HJC) midpoint, while the midpoint itself moved
down 78.7~mm: COM 163.7~mm and L5/S1 217.3~mm above the midpoint against 84.5~mm and
138.2~mm by design. The hip lever arms were therefore wrong.

\subsection{Operation (a): Pure Re-Expression}
Let $F$ be the pelvis body frame of the specification and $F'=FT$ the new frame, with
$T=(R,t)\in SE(3)$ its pose in $F$. The MuJoCo body frame $J$ (the frame the six root
coordinates act on) is not moved, so coordinates keep their meaning. Every quantity
stored in $F$ is re-expressed with $T^{-1}$:
\begin{align}
  p' &= R^\top (p - t), &
  c' &= R^\top (c - t),\quad I'_{c} = R^\top I_{c} R, \\
  P' &= T^{-1} P \text{ (child joint frames, named frames, closure placements)}, &
  C' &= T^{-1} C \text{ (pelvis to $J$, \texttt{child\_to\_follower})},
\end{align}
with $c$ the solid COM and $I_c$ its inertia tensor about $c$. Visual ellipsoids are
axis aligned, so a rotation is refused when the body carries one (translation only).
Markers attach to named frames and are untouched.

\paragraph{Invariance sketch.} A body's world pose is $X_J\,C^{-1}$-composed with its
content: a point $p$ in $F$ is at $X_J C^{-1} p$ in the world ($C$ the pelvis-to-$J$
transform). After the change it is $X_J C'^{-1} p' = X_J (T^{-1}C)^{-1} T^{-1}p
= X_J C^{-1} T T^{-1} p = X_J C^{-1}p$. A child joint frame obeys the same identity, so
child poses are unchanged for every coordinate vector. Mass and the world inertia
$R_w I_c R_w^\top$ are unchanged because $R_w$ absorbs $R$ and $I'_c=R^\top I_c R$.
All terms of the equations of motion are functions of these world quantities and the
unchanged coordinates, so $\mathbf q$, $\dot{\mathbf q}$, $\ddot{\mathbf q}$ and the
generalised forces (\texttt{qfrc\_inverse}) are identical. Composition holds,
$\mathrm{reexpress}(T_1)\circ\mathrm{reexpress}(T_2)=\mathrm{reexpress}(T_1T_2)$, and
$T^{-1}$ undoes $T$ to $10^{-12}$.

\paragraph{Test.} \texttt{tests/unit/motion\_matching/test\_pelvis\_frame.py}: world
landmarks (frame and contact sites), per-body COM and mass, world inertia tensor,
whole-body COM and \texttt{qfrc\_inverse}, relative to the largest entry, $\le10^{-9}$,
for a translation and a $12^\circ,-25^\circ,140^\circ$ rotation plus translation, on
\texttt{full\_body\_spec\_anthro\_driver} and \texttt{full\_body\_spec\_v1}. A mutation
test (tensor left unrotated) fails the same check, so the check is not vacuous. On the
best-fit specification (private, by path): relative differences $0$ for the translation
and $\le 10^{-15}$ (COM $2.9\times10^{-16}$, world inertia $1.0\times10^{-15}$,
\texttt{qfrc\_inverse} $8.8\times10^{-16}$, sites $4.4\times10^{-16}$~m) for the
rotation, with $\max|\mathtt{qfrc}|=615$~N\,m at a random state.

\subsection{Operation (b): Documented Physical Relocation}
Inputs: HJCs $c_R,c_L$ in the \texttt{Hip} (waist marker) frame from HRG-C2
(\texttt{hip\_centre\_map}); status must be \texttt{identified} (fails closed on
\texttt{prior\_dominated} or \texttt{degenerate} unless accepted explicitly).
With $m=\tfrac12(c_R+c_L)$:
\begin{enumerate}
  \item The pelvis frame origin moves to $m$ by operation (a), axes unchanged.
  \item Hip joint translations are set to $c_R,c_L$ ($R$ blocks, hence hip axes and
  zero-twist calibration, unchanged).
  \item The pelvis COM, the L5/S1 pivots (\texttt{TorsoInput} and the spine universal
  joint) and the pelvis visual hint keep their \emph{design offsets from $m$}
  (read from the input; $d_{COM}=(0,0,0.0845)$, $d_{L5S1}=(0,0,0.1382)$~m here).
  \item The bodies hung on L5/S1 keep their fitted pose at zero twist and spine
  angles: the joint pivot moves by $\delta=m$, the follower offset takes the same shift
  ($C_u'=F_u^{-1}\,\mathrm{Trans}(\delta)\,J_u$, likewise for the spine joint). The
  twist and spine rotation centres therefore change, which is the intended effect on
  the L5/S1 torques; trunk, arms, head, markers and total mass do not move.
\end{enumerate}
The output is a NEW document (\texttt{spec\_revision} block: issue, parent canonical
hash, centres, offsets, report) with its own \texttt{spec\_sha256}; the writer refuses
to overwrite a spec or receipt. The parent and all historical receipts are untouched.

\subsection{Independent Pelvic-Plane Reference}
Aligning the pelvis $y$ axis to the hip line makes the two HJC heights equal by
construction, so equal heights then say nothing. The reference here is the waist
marker plane (the \texttt{Hip} frame $z$ axis), and the hip line is a measured quantity
with an interval (\texttt{pelvic\_plane\_report}). On the best capture the MAP height
difference is $-20.2$ mm [$-60.2$, 20.4] and the roll $-6.2^\circ$ [$-18.2$, 6.0]: the
data do not exclude a level hip line, so the frame is \emph{not} rotated to the hip
line. Caveat: the same four markers define the plane and feed the estimate, so common-mode
marker error (13.3 mm waist residual) is shared; there is no ASIS reference in this
capture.

\subsection{Results on the Best Driver Fit}
HJC displacements (original pelvis frame): right $(0.2,\,5.6,\,11.2)$ mm, left
$(2.9,\,-25.4,\,70.6)$ mm; inter-HJC 220.3 to 186.9 mm. L5/S1 to HJC midpoint 217.3 to
138.2 mm (inside 0.13--0.15 m; this equals the design offset by construction, and the
compiled-model geometry confirms it independently). Pelvis COM to midpoint 163.7 to
84.5 mm.

\begin{center}
\begin{tabular}{lcc}
\hline
Address lever arms (horizontal, m, $x$ forward, $y$ left) & Before & After \\
\hline
Pelvis COM from right HJC & $(-0.068,-0.106)$ & $(-0.037,-0.095)$ \\
Pelvis COM from left HJC & $(-0.104,\ 0.109)$ & $(-0.041,\ 0.091)$ \\
Above-hip COM from right HJC & $(-0.232,-0.104)$ & $(-0.226,-0.098)$ \\
Above-hip COM from left HJC & $(-0.269,\ 0.111)$ & $(-0.231,\ 0.088)$ \\
Above-hip COM height over right / left HJC & 0.297 / 0.328 & 0.279 / 0.258 \\
\hline
\end{tabular}
\end{center}

\paragraph{Inverse dynamics.} 91 frames (every 20th of the 1815-frame bundle, 1 ms
sampling, acceleration by central difference of the bundle velocity; impact near frame
1323). The hip reaction wrench is the force and moment about the HJC midpoint that the
legs must apply to the pelvis, from \texttt{qfrc\_inverse} through the root Jacobian,
with the legs given negligible mass (so no leg motion or ground contact is needed). The
L5/S1 torques depend only on the bodies above.
\begin{center}
\begin{tabular}{lcc}
\hline
Peak over the slice & Before & After \\
\hline
Hip reaction force (N) & 1433.5 & 1396.3 \\
Hip reaction moment about HJC midpoint (N\,m) & 316.2 & 414.2 \\
$|\tau|$ \texttt{SpineInputX} (N\,m) & 177.1 & 178.8 \\
$|\tau|$ \texttt{SpineInputY} (N\,m) & 256.2 & 355.2 \\
$|\tau|$ \texttt{TorsoInput} (N\,m) & 222.4 & 244.7 \\
\hline
\end{tabular}
\end{center}
The changes follow from the geometry: the above-hip COM, the pelvis COM and the L5/S1
pivots changed distance to the pivots while the upper-body motion did not. They do not
mean the new torques are more accurate. Stills:
\texttt{HRG\_2026-10-17/HRG-C3/before\_after\_hip\_reaction\_l5s1.png}.

\subsection{Limitations and Failed Attempts}
\begin{itemize}
  \item The left HJC rises 70.6 mm relative to the pelvis markers, so at the fitted
  joint angles the left foot is lifted by the same amount; a contact-aware whole-body
  inverse dynamics on this trajectory would mostly measure that, so none is reported.
  The refit is a later wave.
  \item A first test compared upper-body poses at random spine angles and failed:
  the bodies keep their pose only at zero twist and spine angles, because the rotation
  centre moved. The test now zeroes those coordinates and says so.
  \item The native \texttt{full\_body\_spec\_v1} has spine pivots not colocated with
  the twist pivot, so (b) refuses it (\texttt{ValueError}); (a) works on it.
  \item The root coordinates stay at the body frame $J$; moving them to $m$ needs a state
  map (Euler recovery, velocity, acceleration) and belongs to \texttt{full-body-v2}.
  \item Only a pure translation is applied in (b) and rotations are refused with visual
  ellipsoids; pelvis width is HRG-C4.
\end{itemize}

\subsection{Reproduction}
\begin{verbatim}
# tests (red before the module existed: ImportError at collection)
PYTHONPATH=.:src python3 -m pytest -o addopts="" -q -n 4 \
    tests/unit/motion_matching/test_pelvis_frame.py
# best-fit revision and comparison (private inputs by path; outputs outside the repository)
R=$HOME/fleet24/brick-sim-s32/ct/s32-fw-drv/runs
PYTHONPATH=.:src python3 -m scripts.research.hrg_c3_pelvis_revision \
    --spec $R/driver/full_body_spec_hipcal_scaled.json --receipt $R/driver/receipt.json \
    --capture data/C3D_TA_Driver.c3d --bundle $R/bundles/driver.npz \
    --out-dir $HOME/fleet24/hrg-c3 --stride 20
\end{verbatim}
Revision \texttt{spec\_sha256} 0d08a23a\ldots{} (canonical 314a7f1d\ldots{}) from parent
\texttt{d1829abf\ldots{}}.
.
\section{Pelvis Frame Re-Expression and Hip Centre Relocation}
\label{sec:hrg-pelvis-frame}

\paragraph{Status.} Software correctness only; scientific qualification stays in the
design-manual pathway. No refit was run: the joint trajectory is the fitted bundle's,
so the numbers below describe the model change, not a better match. The fit's inverse
dynamics are replayed state by state (inverse kinematics plus \texttt{mj\_inverse}),
not an independent open-loop simulation.

\subsection{Problem}
\texttt{apply\_hip\_calibration} moves only the hip joints. On the best driver fit
(\texttt{full-body-v1}, 44 coordinates, AN-1 upper body) the pelvis centre of mass
(COM) and L5/S1 stayed at their design offsets from the pelvis origin, which the
builder places at the hip-joint-centre (HJC) midpoint, while the midpoint itself moved
down 78.7~mm: COM 163.7~mm and L5/S1 217.3~mm above the midpoint against 84.5~mm and
138.2~mm by design. The hip lever arms were therefore wrong.

\subsection{Operation (a): Pure Re-Expression}
Let $F$ be the pelvis body frame of the specification and $F'=FT$ the new frame, with
$T=(R,t)\in SE(3)$ its pose in $F$. The MuJoCo body frame $J$ (the frame the six root
coordinates act on) is not moved, so coordinates keep their meaning. Every quantity
stored in $F$ is re-expressed with $T^{-1}$:
\begin{align}
  p' &= R^\top (p - t), &
  c' &= R^\top (c - t),\quad I'_{c} = R^\top I_{c} R, \\
  P' &= T^{-1} P \text{ (child joint frames, named frames, closure placements)}, &
  C' &= T^{-1} C \text{ (pelvis to $J$, \texttt{child\_to\_follower})},
\end{align}
with $c$ the solid COM and $I_c$ its inertia tensor about $c$. Visual ellipsoids are
axis aligned, so a rotation is refused when the body carries one (translation only).
Markers attach to named frames and are untouched.

\paragraph{Invariance sketch.} A body's world pose is $X_J\,C^{-1}$-composed with its
content: a point $p$ in $F$ is at $X_J C^{-1} p$ in the world ($C$ the pelvis-to-$J$
transform). After the change it is $X_J C'^{-1} p' = X_J (T^{-1}C)^{-1} T^{-1}p
= X_J C^{-1} T T^{-1} p = X_J C^{-1}p$. A child joint frame obeys the same identity, so
child poses are unchanged for every coordinate vector. Mass and the world inertia
$R_w I_c R_w^\top$ are unchanged because $R_w$ absorbs $R$ and $I'_c=R^\top I_c R$.
All terms of the equations of motion are functions of these world quantities and the
unchanged coordinates, so $\mathbf q$, $\dot{\mathbf q}$, $\ddot{\mathbf q}$ and the
generalised forces (\texttt{qfrc\_inverse}) are identical. Composition holds,
$\mathrm{reexpress}(T_1)\circ\mathrm{reexpress}(T_2)=\mathrm{reexpress}(T_1T_2)$, and
$T^{-1}$ undoes $T$ to $10^{-12}$.

\paragraph{Test.} \texttt{tests/unit/motion\_matching/test\_pelvis\_frame.py}: world
landmarks (frame and contact sites), per-body COM and mass, world inertia tensor,
whole-body COM and \texttt{qfrc\_inverse}, relative to the largest entry, $\le10^{-9}$,
for a translation and a $12^\circ,-25^\circ,140^\circ$ rotation plus translation, on
\texttt{full\_body\_spec\_anthro\_driver} and \texttt{full\_body\_spec\_v1}. A mutation
test (tensor left unrotated) fails the same check, so the check is not vacuous. On the
best-fit specification (private, by path): relative differences $0$ for the translation
and $\le 10^{-15}$ (COM $2.9\times10^{-16}$, world inertia $1.0\times10^{-15}$,
\texttt{qfrc\_inverse} $8.8\times10^{-16}$, sites $4.4\times10^{-16}$~m) for the
rotation, with $\max|\mathtt{qfrc}|=615$~N\,m at a random state.

\subsection{Operation (b): Documented Physical Relocation}
Inputs: HJCs $c_R,c_L$ in the \texttt{Hip} (waist marker) frame from HRG-C2
(\texttt{hip\_centre\_map}); status must be \texttt{identified} (fails closed on
\texttt{prior\_dominated} or \texttt{degenerate} unless accepted explicitly).
With $m=\tfrac12(c_R+c_L)$:
\begin{enumerate}
  \item The pelvis frame origin moves to $m$ by operation (a), axes unchanged.
  \item Hip joint translations are set to $c_R,c_L$ ($R$ blocks, hence hip axes and
  zero-twist calibration, unchanged).
  \item The pelvis COM, the L5/S1 pivots (\texttt{TorsoInput} and the spine universal
  joint) and the pelvis visual hint keep their \emph{design offsets from $m$}
  (read from the input; $d_{COM}=(0,0,0.0845)$, $d_{L5S1}=(0,0,0.1382)$~m here).
  \item The bodies hung on L5/S1 keep their fitted pose at zero twist and spine
  angles: the joint pivot moves by $\delta=m$, the follower offset takes the same shift
  ($C_u'=F_u^{-1}\,\mathrm{Trans}(\delta)\,J_u$, likewise for the spine joint). The
  twist and spine rotation centres therefore change, which is the intended effect on
  the L5/S1 torques; trunk, arms, head, markers and total mass do not move.
\end{enumerate}
The output is a NEW document (\texttt{spec\_revision} block: issue, parent canonical
hash, centres, offsets, report) with its own \texttt{spec\_sha256}; the writer refuses
to overwrite a spec or receipt. The parent and all historical receipts are untouched.

\subsection{Independent Pelvic-Plane Reference}
Aligning the pelvis $y$ axis to the hip line makes the two HJC heights equal by
construction, so equal heights then say nothing. The reference here is the waist
marker plane (the \texttt{Hip} frame $z$ axis), and the hip line is a measured quantity
with an interval (\texttt{pelvic\_plane\_report}). On the best capture the MAP height
difference is $-20.2$ mm [$-60.2$, 20.4] and the roll $-6.2^\circ$ [$-18.2$, 6.0]: the
data do not exclude a level hip line, so the frame is \emph{not} rotated to the hip
line. Caveat: the same four markers define the plane and feed the estimate, so common-mode
marker error (13.3 mm waist residual) is shared; there is no ASIS reference in this
capture.

\subsection{Results on the Best Driver Fit}
HJC displacements (original pelvis frame): right $(0.2,\,5.6,\,11.2)$ mm, left
$(2.9,\,-25.4,\,70.6)$ mm; inter-HJC 220.3 to 186.9 mm. L5/S1 to HJC midpoint 217.3 to
138.2 mm (inside 0.13--0.15 m; this equals the design offset by construction, and the
compiled-model geometry confirms it independently). Pelvis COM to midpoint 163.7 to
84.5 mm.

\begin{center}
\begin{tabular}{lcc}
\hline
Address lever arms (horizontal, m, $x$ forward, $y$ left) & Before & After \\
\hline
Pelvis COM from right HJC & $(-0.068,-0.106)$ & $(-0.037,-0.095)$ \\
Pelvis COM from left HJC & $(-0.104,\ 0.109)$ & $(-0.041,\ 0.091)$ \\
Above-hip COM from right HJC & $(-0.232,-0.104)$ & $(-0.226,-0.098)$ \\
Above-hip COM from left HJC & $(-0.269,\ 0.111)$ & $(-0.231,\ 0.088)$ \\
Above-hip COM height over right / left HJC & 0.297 / 0.328 & 0.279 / 0.258 \\
\hline
\end{tabular}
\end{center}

\paragraph{Inverse dynamics.} 91 frames (every 20th of the 1815-frame bundle, 1 ms
sampling, acceleration by central difference of the bundle velocity; impact near frame
1323). The hip reaction wrench is the force and moment about the HJC midpoint that the
legs must apply to the pelvis, from \texttt{qfrc\_inverse} through the root Jacobian,
with the legs given negligible mass (so no leg motion or ground contact is needed). The
L5/S1 torques depend only on the bodies above.
\begin{center}
\begin{tabular}{lcc}
\hline
Peak over the slice & Before & After \\
\hline
Hip reaction force (N) & 1433.5 & 1396.3 \\
Hip reaction moment about HJC midpoint (N\,m) & 316.2 & 414.2 \\
$|\tau|$ \texttt{SpineInputX} (N\,m) & 177.1 & 178.8 \\
$|\tau|$ \texttt{SpineInputY} (N\,m) & 256.2 & 355.2 \\
$|\tau|$ \texttt{TorsoInput} (N\,m) & 222.4 & 244.7 \\
\hline
\end{tabular}
\end{center}
The changes follow from the geometry: the above-hip COM, the pelvis COM and the L5/S1
pivots changed distance to the pivots while the upper-body motion did not. They do not
mean the new torques are more accurate. Stills:
\texttt{HRG\_2026-10-17/HRG-C3/before\_after\_hip\_reaction\_l5s1.png}.

\subsection{Limitations and Failed Attempts}
\begin{itemize}
  \item The left HJC rises 70.6 mm relative to the pelvis markers, so at the fitted
  joint angles the left foot is lifted by the same amount; a contact-aware whole-body
  inverse dynamics on this trajectory would mostly measure that, so none is reported.
  The refit is a later wave.
  \item A first test compared upper-body poses at random spine angles and failed:
  the bodies keep their pose only at zero twist and spine angles, because the rotation
  centre moved. The test now zeroes those coordinates and says so.
  \item The native \texttt{full\_body\_spec\_v1} has spine pivots not colocated with
  the twist pivot, so (b) refuses it (\texttt{ValueError}); (a) works on it.
  \item The root coordinates stay at the body frame $J$; moving them to $m$ needs a state
  map (Euler recovery, velocity, acceleration) and belongs to \texttt{full-body-v2}.
  \item Only a pure translation is applied in (b) and rotations are refused with visual
  ellipsoids; pelvis width is HRG-C4.
\end{itemize}

\subsection{Reproduction}
\begin{verbatim}
# tests (red before the module existed: ImportError at collection)
PYTHONPATH=.:src python3 -m pytest -o addopts="" -q -n 4 \
    tests/unit/motion_matching/test_pelvis_frame.py
# best-fit revision and comparison (private inputs by path; outputs outside the repository)
R=$HOME/fleet24/brick-sim-s32/ct/s32-fw-drv/runs
PYTHONPATH=.:src python3 -m scripts.research.hrg_c3_pelvis_revision \
    --spec $R/driver/full_body_spec_hipcal_scaled.json --receipt $R/driver/receipt.json \
    --capture data/C3D_TA_Driver.c3d --bundle $R/bundles/driver.npz \
    --out-dir $HOME/fleet24/hrg-c3 --stride 20
\end{verbatim}
Revision \texttt{spec\_sha256} 0d08a23a\ldots{} (canonical 314a7f1d\ldots{}) from parent
\texttt{d1829abf\ldots{}}.
.
\section{Pelvis Frame Re-Expression and Hip Centre Relocation}
\label{sec:hrg-pelvis-frame}

\paragraph{Status.} Software correctness only; scientific qualification stays in the
design-manual pathway. No refit was run: the joint trajectory is the fitted bundle's,
so the numbers below describe the model change, not a better match. The fit's inverse
dynamics are replayed state by state (inverse kinematics plus \texttt{mj\_inverse}),
not an independent open-loop simulation.

\subsection{Problem}
\texttt{apply\_hip\_calibration} moves only the hip joints. On the best driver fit
(\texttt{full-body-v1}, 44 coordinates, AN-1 upper body) the pelvis centre of mass
(COM) and L5/S1 stayed at their design offsets from the pelvis origin, which the
builder places at the hip-joint-centre (HJC) midpoint, while the midpoint itself moved
down 78.7~mm: COM 163.7~mm and L5/S1 217.3~mm above the midpoint against 84.5~mm and
138.2~mm by design. The hip lever arms were therefore wrong.

\subsection{Operation (a): Pure Re-Expression}
Let $F$ be the pelvis body frame of the specification and $F'=FT$ the new frame, with
$T=(R,t)\in SE(3)$ its pose in $F$. The MuJoCo body frame $J$ (the frame the six root
coordinates act on) is not moved, so coordinates keep their meaning. Every quantity
stored in $F$ is re-expressed with $T^{-1}$:
\begin{align}
  p' &= R^\top (p - t), &
  c' &= R^\top (c - t),\quad I'_{c} = R^\top I_{c} R, \\
  P' &= T^{-1} P \text{ (child joint frames, named frames, closure placements)}, &
  C' &= T^{-1} C \text{ (pelvis to $J$, \texttt{child\_to\_follower})},
\end{align}
with $c$ the solid COM and $I_c$ its inertia tensor about $c$. Visual ellipsoids are
axis aligned, so a rotation is refused when the body carries one (translation only).
Markers attach to named frames and are untouched.

\paragraph{Invariance sketch.} A body's world pose is $X_J\,C^{-1}$-composed with its
content: a point $p$ in $F$ is at $X_J C^{-1} p$ in the world ($C$ the pelvis-to-$J$
transform). After the change it is $X_J C'^{-1} p' = X_J (T^{-1}C)^{-1} T^{-1}p
= X_J C^{-1} T T^{-1} p = X_J C^{-1}p$. A child joint frame obeys the same identity, so
child poses are unchanged for every coordinate vector. Mass and the world inertia
$R_w I_c R_w^\top$ are unchanged because $R_w$ absorbs $R$ and $I'_c=R^\top I_c R$.
All terms of the equations of motion are functions of these world quantities and the
unchanged coordinates, so $\mathbf q$, $\dot{\mathbf q}$, $\ddot{\mathbf q}$ and the
generalised forces (\texttt{qfrc\_inverse}) are identical. Composition holds,
$\mathrm{reexpress}(T_1)\circ\mathrm{reexpress}(T_2)=\mathrm{reexpress}(T_1T_2)$, and
$T^{-1}$ undoes $T$ to $10^{-12}$.

\paragraph{Test.} \texttt{tests/unit/motion\_matching/test\_pelvis\_frame.py}: world
landmarks (frame and contact sites), per-body COM and mass, world inertia tensor,
whole-body COM and \texttt{qfrc\_inverse}, relative to the largest entry, $\le10^{-9}$,
for a translation and a $12^\circ,-25^\circ,140^\circ$ rotation plus translation, on
\texttt{full\_body\_spec\_anthro\_driver} and \texttt{full\_body\_spec\_v1}. A mutation
test (tensor left unrotated) fails the same check, so the check is not vacuous. On the
best-fit specification (private, by path): relative differences $0$ for the translation
and $\le 10^{-15}$ (COM $2.9\times10^{-16}$, world inertia $1.0\times10^{-15}$,
\texttt{qfrc\_inverse} $8.8\times10^{-16}$, sites $4.4\times10^{-16}$~m) for the
rotation, with $\max|\mathtt{qfrc}|=615$~N\,m at a random state.

\subsection{Operation (b): Documented Physical Relocation}
Inputs: HJCs $c_R,c_L$ in the \texttt{Hip} (waist marker) frame from HRG-C2
(\texttt{hip\_centre\_map}); status must be \texttt{identified} (fails closed on
\texttt{prior\_dominated} or \texttt{degenerate} unless accepted explicitly).
With $m=\tfrac12(c_R+c_L)$:
\begin{enumerate}
  \item The pelvis frame origin moves to $m$ by operation (a), axes unchanged.
  \item Hip joint translations are set to $c_R,c_L$ ($R$ blocks, hence hip axes and
  zero-twist calibration, unchanged).
  \item The pelvis COM, the L5/S1 pivots (\texttt{TorsoInput} and the spine universal
  joint) and the pelvis visual hint keep their \emph{design offsets from $m$}
  (read from the input; $d_{COM}=(0,0,0.0845)$, $d_{L5S1}=(0,0,0.1382)$~m here).
  \item The bodies hung on L5/S1 keep their fitted pose at zero twist and spine
  angles: the joint pivot moves by $\delta=m$, the follower offset takes the same shift
  ($C_u'=F_u^{-1}\,\mathrm{Trans}(\delta)\,J_u$, likewise for the spine joint). The
  twist and spine rotation centres therefore change, which is the intended effect on
  the L5/S1 torques; trunk, arms, head, markers and total mass do not move.
\end{enumerate}
The output is a NEW document (\texttt{spec\_revision} block: issue, parent canonical
hash, centres, offsets, report) with its own \texttt{spec\_sha256}; the writer refuses
to overwrite a spec or receipt. The parent and all historical receipts are untouched.

\subsection{Independent Pelvic-Plane Reference}
Aligning the pelvis $y$ axis to the hip line makes the two HJC heights equal by
construction, so equal heights then say nothing. The reference here is the waist
marker plane (the \texttt{Hip} frame $z$ axis), and the hip line is a measured quantity
with an interval (\texttt{pelvic\_plane\_report}). On the best capture the MAP height
difference is $-20.2$ mm [$-60.2$, 20.4] and the roll $-6.2^\circ$ [$-18.2$, 6.0]: the
data do not exclude a level hip line, so the frame is \emph{not} rotated to the hip
line. Caveat: the same four markers define the plane and feed the estimate, so common-mode
marker error (13.3 mm waist residual) is shared; there is no ASIS reference in this
capture.

\subsection{Results on the Best Driver Fit}
HJC displacements (original pelvis frame): right $(0.2,\,5.6,\,11.2)$ mm, left
$(2.9,\,-25.4,\,70.6)$ mm; inter-HJC 220.3 to 186.9 mm. L5/S1 to HJC midpoint 217.3 to
138.2 mm (inside 0.13--0.15 m; this equals the design offset by construction, and the
compiled-model geometry confirms it independently). Pelvis COM to midpoint 163.7 to
84.5 mm.

\begin{center}
\begin{tabular}{lcc}
\hline
Address lever arms (horizontal, m, $x$ forward, $y$ left) & Before & After \\
\hline
Pelvis COM from right HJC & $(-0.068,-0.106)$ & $(-0.037,-0.095)$ \\
Pelvis COM from left HJC & $(-0.104,\ 0.109)$ & $(-0.041,\ 0.091)$ \\
Above-hip COM from right HJC & $(-0.232,-0.104)$ & $(-0.226,-0.098)$ \\
Above-hip COM from left HJC & $(-0.269,\ 0.111)$ & $(-0.231,\ 0.088)$ \\
Above-hip COM height over right / left HJC & 0.297 / 0.328 & 0.279 / 0.258 \\
\hline
\end{tabular}
\end{center}

\paragraph{Inverse dynamics.} 91 frames (every 20th of the 1815-frame bundle, 1 ms
sampling, acceleration by central difference of the bundle velocity; impact near frame
1323). The hip reaction wrench is the force and moment about the HJC midpoint that the
legs must apply to the pelvis, from \texttt{qfrc\_inverse} through the root Jacobian,
with the legs given negligible mass (so no leg motion or ground contact is needed). The
L5/S1 torques depend only on the bodies above.
\begin{center}
\begin{tabular}{lcc}
\hline
Peak over the slice & Before & After \\
\hline
Hip reaction force (N) & 1433.5 & 1396.3 \\
Hip reaction moment about HJC midpoint (N\,m) & 316.2 & 414.2 \\
$|\tau|$ \texttt{SpineInputX} (N\,m) & 177.1 & 178.8 \\
$|\tau|$ \texttt{SpineInputY} (N\,m) & 256.2 & 355.2 \\
$|\tau|$ \texttt{TorsoInput} (N\,m) & 222.4 & 244.7 \\
\hline
\end{tabular}
\end{center}
The changes follow from the geometry: the above-hip COM, the pelvis COM and the L5/S1
pivots changed distance to the pivots while the upper-body motion did not. They do not
mean the new torques are more accurate. Stills:
\texttt{HRG\_2026-10-17/HRG-C3/before\_after\_hip\_reaction\_l5s1.png}.

\subsection{Limitations and Failed Attempts}
\begin{itemize}
  \item The left HJC rises 70.6 mm relative to the pelvis markers, so at the fitted
  joint angles the left foot is lifted by the same amount; a contact-aware whole-body
  inverse dynamics on this trajectory would mostly measure that, so none is reported.
  The refit is a later wave.
  \item A first test compared upper-body poses at random spine angles and failed:
  the bodies keep their pose only at zero twist and spine angles, because the rotation
  centre moved. The test now zeroes those coordinates and says so.
  \item The native \texttt{full\_body\_spec\_v1} has spine pivots not colocated with
  the twist pivot, so (b) refuses it (\texttt{ValueError}); (a) works on it.
  \item The root coordinates stay at the body frame $J$; moving them to $m$ needs a state
  map (Euler recovery, velocity, acceleration) and belongs to \texttt{full-body-v2}.
  \item Only a pure translation is applied in (b) and rotations are refused with visual
  ellipsoids; pelvis width is HRG-C4.
\end{itemize}

\subsection{Reproduction}
\begin{verbatim}
# tests (red before the module existed: ImportError at collection)
PYTHONPATH=.:src python3 -m pytest -o addopts="" -q -n 4 \
    tests/unit/motion_matching/test_pelvis_frame.py
# best-fit revision and comparison (private inputs by path; outputs outside the repository)
R=$HOME/fleet24/brick-sim-s32/ct/s32-fw-drv/runs
PYTHONPATH=.:src python3 -m scripts.research.hrg_c3_pelvis_revision \
    --spec $R/driver/full_body_spec_hipcal_scaled.json --receipt $R/driver/receipt.json \
    --capture data/C3D_TA_Driver.c3d --bundle $R/bundles/driver.npz \
    --out-dir $HOME/fleet24/hrg-c3 --stride 20
\end{verbatim}
Revision \texttt{spec\_sha256} 0d08a23a\ldots{} (canonical 314a7f1d\ldots{}) from parent
\texttt{d1829abf\ldots{}}.
